\subsection{Data Collection}

Our data come from a single Unitree Go2 quadruped, whose onboard proprioception is published on the ROS~2 topic \texttt{/lowstate}~\cite{macenski2022ros2,unitree_go2}. We recorded this stream while running a set of scripted locomotion scenarios, each designed to place the robot in a controlled situation and to log everything the joints, feet, and body sensors reported during it. Figure~\ref{fig:data} shows the collection pipeline end to end, from the live \texttt{/lowstate} messages through per-recording comma-separated-value (CSV) logs to the manually labeled dataset used for training and evaluation.

Each scripted scenario walks the robot for a while and, in the positive cases, drives one leg into a physical entanglement (for example, a wire or a hand snagging the limb) before releasing it. We also scripted purely negative scenarios that contain no entanglement at all, as well as a deliberate \texttt{Lock} condition: a rigid stance the robot holds immediately after standing up. \texttt{Lock} matters because it produces large, sustained joint torques with almost no motion, which superficially resembles a snag; including it as a scripted negative lets the model learn to tell a held stance apart from a genuine entanglement.

Logging is event-driven rather than clocked, so the effective sample rate is not uniform. Across recordings it ranges from roughly $410$ to $500$~Hz, with a median near $483$~Hz. We keep this raw, non-uniform timing in the recorded logs and defer resampling to the feature stage, so that no timing assumption is baked into the collected data.

The full corpus is $25$ recordings totaling $159{,}001$ labeled rows. Each row has $51$ columns, which fall into three sensor groups plus a label. The motor group holds $36$ values: the four legs, each with a hip, thigh, and calf joint, and for every joint a position $q$, a velocity $dq$, and a torque $\tau$. The next group is the four foot-contact forces, one per leg. The last sensor group is the inertial measurement unit (IMU), reporting body roll, pitch, and yaw, the three gyroscope axes, and the three accelerometer axes. A final \texttt{Status} column carries the phase label. Table~\ref{tab:dataset} summarizes these groups and the recording counts.

We labeled the recordings by hand. Because every scenario was scripted, we knew when an entanglement was induced and released, and we used those recorded timestamps together with per-joint torque plots to mark the exact rows spanning each event. The torque traces are the more reliable cue: a snagged limb shows a distinctive sustained resistive torque, which is visible in the plots even when the wall-clock timestamps drift slightly from the logged samples. Each row receives one of the phase labels \texttt{Walking}, \texttt{Entangled}, \texttt{Stop}, \texttt{Lock}, or blank, and the affected leg for a positive recording is encoded in that recording's filename.

Of the $25$ recordings, $15$ are positive (each containing one contiguous entangled event) and $10$ are negative. The positive events are not evenly distributed across the four legs. Grouping by the affected leg gives seven recordings for the rear-right leg, five for the rear-left, four for the front-left, and only three for the front-right. The front-right leg is therefore under-represented, and we flag it here because it constrains how confidently the per-leg behavior for that limb can be estimated; the effect resurfaces in the per-leg results (Figure~\ref{fig:perleg}). Figure~\ref{fig:dataset} plots the distribution of labels and per-leg positives across the dataset.

The way recordings are split into training, validation, and test sets is central to the credibility of every number we report, so we describe it before any modeling detail. We split at the level of whole recordings, assigning $16$ recordings to training, $4$ to validation, and $5$ to test, so that no recording contributes rows to more than one partition. The alternative, cutting the pooled rows or sliding windows at random, would be a mistake: consecutive windows from one recording overlap heavily and share the same robot, the same scripted scenario, and the same entanglement event, so a window in the test set would be nearly identical to windows the model trained on. That kind of window-level split leaks the answer and inflates the metrics. Grouping by recording forces the model to generalize to events it has never seen. The held-out test set is \texttt{back\_left\_hand2}, \texttt{back\_right\_wire1}, \texttt{front\_left\_hand2}, and \texttt{front\_both\_wire2} as positives, with \texttt{go2\_lowstate\_lock\_stop} as the negative, giving four positive recordings and one negative.